% Compilar desde la raíz del repositorio.
% Versión v5: se conserva el temario y las 13 imágenes del tema original.
% Se añade contexto didáctico para alumnado de 1.º sin ampliar innecesariamente
% el número de diapositivas: 24 frente a las 22 originales.
% Cambios principales:
% - recordatorio breve de dominio e imagen;
% - conexión explícita entre funciones de una y de varias variables;
% - resumen de restricciones habituales para calcular dominios;
% - explicación intuitiva de las curvas de nivel como cortes a altura constante;
% - contexto breve en las diapositivas que eran solo imágenes;
% - separación de los ejercicios finales de dominio y conjuntos de nivel.
\input{config.tex}
\begin{document}

% Diapositiva 1
\begin{frame}{}
\centering
\LARGE Funciones de varias variables
\par\vspace{1cm}
\large Grado en Ingeniería Informática
\par Matemáticas 2
\end{frame}

% Diapositiva 2
\begin{frame}{Ejemplo: temperatura}
La temperatura $T$ en un punto de la superficie de la Tierra a una hora determinada depende de la longitud y la latitud de dicho punto.
\medskip

Se puede pensar en $T$ como una función de dos variables y, si las denominamos $x$ e $y$, esta dependencia funcional se puede escribir como
\[
T=f(x,y).
\]
\smallskip

La entrada ya no es un único número, sino un par $(x,y)$ que indica una posición. La salida sigue siendo un único valor: la temperatura.
\end{frame}

% Diapositiva 3
\begin{frame}{Ejemplo: volumen}
El volumen de una caja rectangular de dimensiones $x,y,z$ es
\[
\text{Volumen}=x\cdot y\cdot z.
\]
Este es un ejemplo de una función real de tres variables reales:
\[
f(x,y,z)=x\cdot y\cdot z.
\]
Por ejemplo,
\[
f(2,3,4)=24.
\]
Si las dimensiones están en metros, el resultado es un volumen de $24\,\mathrm{m}^3$.
\figura{s3_2.png}{.40}{.31}
\end{frame}

% Diapositiva 4
\begin{frame}{Definición general}
Una $n$-tupla es una lista ordenada de $n$ números: un par contiene dos, una terna contiene tres, etc.
\medskip

\textbf{Definición}
\medskip

En general, una función real de $n$ variables reales es una correspondencia que asigna a cada $n$-tupla
\[
(x_1,x_2,\ldots,x_n)
\]
de su dominio un único valor
\[
u=f(x_1,x_2,\ldots,x_n).
\]
\end{frame}

% Diapositiva 5
\begin{frame}{Dominio e imagen}
Los conceptos conocidos para funciones de una variable tienen su equivalente para funciones de varias variables.
\medskip

Por ejemplo, para
\[
f(x)=x^2,
\]
el \textbf{dominio} es el conjunto de valores que podemos introducir y la \textbf{imagen} es el conjunto de valores que la función alcanza:
\[
\operatorname{Dom}(f)=\mathbb{R},
\qquad
\operatorname{Im}(f)=[0,\infty).
\]
\medskip

En una variable la entrada es un número; en dos variables, la entrada es un par $(x,y)$; en tres variables, una terna $(x,y,z)$.
\end{frame}

% Diapositiva 6
\begin{frame}{Dos variables}
Particularizando a dos variables, una función $f$ asigna a cada par ordenado de números reales $(x,y)$ perteneciente a un conjunto $D$ un único número real, denotado por
\[
u=f(x,y).
\]
El conjunto $D$ es el dominio de $f$ y su imagen es el conjunto de valores que toma $f$.
\medskip

La idea es la misma que en una función de una variable: damos una entrada y obtenemos una salida. La diferencia es que ahora la entrada contiene dos números.
\figura{s6_0.png}{.86}{.27}
\end{frame}

% Diapositiva 7
\begin{frame}{Notación}
Las funciones de dos variables se suelen escribir de la forma
\[
z=f(x,y),
\]
para explicitar el valor que toma la función en un punto $(x,y)$.
\medskip

Las variables $x$ e $y$ son las \textbf{variables independientes} y $z$ es la \textbf{variable dependiente}.
\medskip

Por ejemplo, si
\[
z=x+y,
\]
entonces la entrada $(2,3)$ produce la salida $z=5$.
\end{frame}

% Diapositiva 8
\begin{frame}{Representación gráfica}
\small
Si $f$ es una función de dos variables con dominio $D$, su gráfica es el conjunto de ternas $(x,y,z)$ de $\mathbb{R}^3$ tales que
\[
z=f(x,y),\qquad (x,y)\in D.
\]
\medskip

En una función de una variable, cada entrada $x$ produce un punto $(x,f(x))$, por lo que la gráfica es una curva en el plano.
\medskip

En una función de dos variables, cada entrada $(x,y)$ produce un punto $(x,y,f(x,y))$, por lo que la gráfica es, en general, una superficie en el espacio.
\figura{s8_2.png}{.40}{.37}
\end{frame}

% Diapositiva 9
\begin{frame}{Ejemplo de gráfica}
Por ejemplo,
\[
z=f(x,y)=x+y^2
\]
es una función de dos variables cuyo dominio es todo $\mathbb{R}^2$.
\medskip

Su representación en 3D es:
\figura{s9_2.png}{.76}{.48}

\small
Podemos entender algunos cortes sencillos: si $x=0$, queda $z=y^2$; si $y=0$, queda $z=x$.
\end{frame}

% Diapositiva 10
\begin{frame}{Restricciones del dominio}
Para encontrar el dominio buscamos qué operaciones pueden dejar de estar definidas.
\medskip

\begin{center}
\begin{tabular}{ll}
\textbf{Expresión} & \textbf{Condición} \\ \hline
$\dfrac{1}{t}$ & $t\neq 0$ \\[2mm]
$\sqrt{t}$ & $t\geq 0$ \\[2mm]
$\ln(t)$ & $t>0$ \\[2mm]
$\arcsin(t)$ & $-1\leq t\leq 1$
\end{tabular}
\end{center}
\medskip

En varias variables se aplican exactamente las mismas reglas, pero ahora $t$ puede ser una expresión que dependa de $x$, $y$, $z$, etc.
\end{frame}

% Diapositiva 11
\begin{frame}{Dominio: cociente}
\textbf{Ejemplo}
\[
z=f(x,y)=\frac{3x-5y}{y-x^2}.
\]
El denominador no puede ser cero:
\[
y-x^2\neq 0.
\]
Por tanto,
\[
D=\{(x,y)\in\mathbb{R}^2:y-x^2\ne0\}.
\]
\figura{s10_4.png}{.66}{.35}

\small
El dominio es todo el plano $\mathbb{R}^2$ excepto los puntos de la parábola $y=x^2$.
\end{frame}

% Diapositiva 12
\begin{frame}{Dominio: raíz}
\textbf{Ejemplo}
\[
z=f(x,y)=\sqrt{9-x^2-y^2}.
\]
El contenido de la raíz debe ser no negativo:
\[
9-x^2-y^2\geq 0
\quad\Longleftrightarrow\quad
x^2+y^2\leq 9.
\]
Así,
\[
D=\{(x,y)\in\mathbb{R}^2:x^2+y^2\leq9\}.
\]
\figura{s11_5.png}{.70}{.32}

\small
El dominio es el interior y la frontera de la circunferencia de centro el origen y radio $3$.
\end{frame}

% Diapositiva 13
\begin{frame}{Dominio: logaritmo}
El dominio de la función
\[
f(x,y,z)=\ln(z-y)+xy\sin z
\]
se obtiene imponiendo que el argumento del logaritmo sea positivo:
\[
z-y>0
\quad\Longleftrightarrow\quad
z>y.
\]
El término $xy\sin z$ no introduce ninguna restricción adicional.
\medskip

Por tanto,
\[
D=\{(x,y,z)\in\mathbb{R}^3:z>y\}.
\]
Es decir, todos los puntos situados por encima del plano $z=y$.
\end{frame}

% Diapositiva 14
\begin{frame}{Operaciones}
Las funciones de varias variables pueden operarse de forma análoga a las de una variable: suma, producto, cociente o composición, siempre que las expresiones resultantes estén definidas.
\medskip

Por ejemplo, si
\[
f(x,y)=x+y,
\qquad
g(x,y)=xy,
\]
entonces
\[
(f+g)(x,y)=x+y+xy,
\]
y
\[
(fg)(x,y)=(x+y)xy.
\]
\smallskip

En un cociente aparecería además la condición de que el denominador no fuese cero.
\end{frame}

% Diapositiva 15
\begin{frame}{Curvas de nivel}
Las curvas de nivel de una función $f$ de dos variables son las curvas definidas por
\[
f(x,y)=k,
\]
donde $k$ es una constante perteneciente a la imagen de $f$.
\medskip

Una curva de nivel reúne todos los puntos $(x,y)$ en los que la función toma el mismo valor.
\medskip

Geométricamente, podemos imaginar que cortamos la superficie $z=f(x,y)$ con un plano horizontal $z=k$ y observamos el corte desde arriba.
\begin{columns}[c]
\begin{column}{.48\textwidth}\figura{s14_3.png}{.80}{.38}\end{column}
\begin{column}{.48\textwidth}\figura{s14_5.png}{.92}{.38}\end{column}
\end{columns}
\end{frame}

% Diapositiva 16
\begin{frame}{Ejemplo: semiesfera}
\textbf{Ejemplo}
\medskip

Obtener las curvas de nivel de
\[
z=\sqrt{100-x^2-y^2}.
\]
\medskip

La gráfica es la semiesfera superior de radio $10$.
\medskip

Su dominio es
\[
x^2+y^2\leq 100,
\]
y sus valores de altura están entre $0$ y $10$.
\medskip

Para obtener una curva de nivel fijaremos una altura constante $z=c$.
\end{frame}

% Diapositiva 17
\begin{frame}{Cálculo de niveles}
\small
Para
\[
z=\sqrt{100-x^2-y^2},
\]
fijamos una altura $z=c$:
\[
c=\sqrt{100-x^2-y^2}.
\]
Elevando al cuadrado,
\[
c^2=100-x^2-y^2,
\]
y reorganizando,
\[
x^2+y^2=100-c^2.
\]
\medskip

Para $0\leq c<10$ obtenemos circunferencias de centro el origen. Para $c=10$ obtenemos únicamente el punto $(0,0)$. Para $c<0$ o $c>10$ no hay puntos en el conjunto de nivel.
\end{frame}

% Diapositiva 18
\begin{frame}{Representación de niveles}
\begin{columns}[c]
\begin{column}{.43\textwidth}
Para
\[
z=\sqrt{100-x^2-y^2},
\]
las curvas de nivel satisfacen
\[
x^2+y^2=100-c^2.
\]
\medskip

Al aumentar la altura $c$, disminuye el radio de la circunferencia.
\end{column}
\begin{column}{.55\textwidth}
\figura{s17_3.png}{1}{.58}
\end{column}
\end{columns}
\end{frame}

% Diapositiva 19
\begin{frame}{Curvas de nivel}
\figura{s18_0.jpg}{.70}{.68}
\begin{center}
\small Cada curva reúne puntos de la superficie en los que la función toma el mismo valor.
\end{center}
\end{frame}

% Diapositiva 20
\begin{frame}{Lectura de niveles}
\figura{s19_1.jpg}{.94}{.58}
\begin{center}
\small El dibujo de las curvas de nivel permite estudiar la forma de una superficie sin representarla directamente en tres dimensiones.
\end{center}
\end{frame}

% Diapositiva 21
\begin{frame}{Separación entre niveles}
\figura{s20_1.jpg}{.94}{.61}
\begin{center}
\small Cuando los niveles representados están igualmente espaciados, curvas más próximas indican un cambio más rápido del valor de la función.
\end{center}
\end{frame}

% Diapositiva 22
\begin{frame}{Mapa topográfico}
\figura{s21_0.jpg}{.66}{.69}
\begin{center}
\small En un mapa topográfico, cada curva une puntos con la misma altura. Es una aplicación directa del concepto de curva de nivel.
\end{center}
\end{frame}

% Diapositiva 23
\begin{frame}{Ejercicio de dominio}
Encontrar el dominio de la función
\[
f(x,y,z)=\arcsin(xy)\,e^{3z}+e^{(y+z)\ln(x+z)}.
\]
\medskip

\textbf{Pista:} identifica primero las operaciones que introducen restricciones:
\[
\arcsin(t):\quad -1\leq t\leq1,
\qquad
\ln(t):\quad t>0.
\]
\medskip

Después aplica cada condición al argumento correspondiente y reúne todas las restricciones.
\end{frame}

% Diapositiva 24
\begin{frame}{Ejercicios de nivel}
Describir los conjuntos de nivel para los valores de $k$ indicados:
\begin{enumerate}[a)]
\item $f(x,y)=x^2+y^2$,\quad $k=0,1,2,3$.
\item $f(x,y,z)=x^2+y^2$,\quad $k=0,1,2$.
\item $f(x,y)=x^2-y^2$,\quad $k=-2,-1,0,1,2$.
\end{enumerate}
\medskip

En a) y c) se obtienen curvas de nivel en el plano. En b), al intervenir tres variables, los conjuntos de nivel son superficies en el espacio.
\medskip

\textbf{Procedimiento:} sustituye $f$ por $k$, simplifica la ecuación e identifica la figura geométrica resultante.
\end{frame}

\end{document}
